\subsection{DECOMPOSITION OF REPRESENTATIONS OF A NILPOTENT LIE ALGEBRA}

Let \(\mathfrak{h}\) be a Lie algebra and M an \(\mathfrak{h}\) -module. For any map \(\lambda\) from \(\mathfrak{h}\) to \(k\) , we have defined in no. 1 vector subspaces \(\mathrm{M}^{\lambda}(\mathfrak{h})\) and \(\mathrm{M}_{\lambda}(\mathfrak{h})\) of M. In particular, if \(\mathfrak{g}\) is a Lie algebra containing \(\mathfrak{h}\) as a subalgebra, and if \(x \in \mathfrak{g}\) , we shall often employ the notations \(\mathfrak{g}^{\lambda}(\mathfrak{h})\) and \(\mathfrak{g}_{\lambda}(\mathfrak{h})\) ; it will then be understood that \(\mathfrak{h}\) operates on \(\mathfrak{g}\) by the adjoint representation ad \(_{\mathfrak{g}}\) .

PROPOSITION 8. Let \(\mathfrak{h}\) be a Lie algebra, and L, M, N \(\mathfrak{h}\) -modules. Denote by P the set of maps from \(\mathfrak{h}\) to \(k\) .

\begin{enumerate}
\def\labelenumi{(\roman{enumi})}
\item
  The sum \(\sum_{\lambda \in \mathrm{P}}\mathrm{L}^{\lambda}(\mathrm{P})\) is direct.
\item
  If \(f: \mathrm{L} \to \mathrm{M}\) is a homomorphism of \(\mathfrak{h}\) -modules, \(f(\mathrm{L}^{\lambda}(\mathfrak{h})) \subset \mathrm{M}^{\lambda}(\mathfrak{h})\) for all \(\lambda \in \mathrm{P}\) .
\item
  If \(f:\mathrm{L}\times \mathrm{M}\to \mathrm{N}\) is a bilinear \(\mathfrak{h}\) -invariant map,
\end{enumerate}

\[
f (\mathrm{L} ^ {\lambda} (\mathfrak {h}) \times \mathrm{M} ^ {\mu} (\mathfrak {h})) \subset \mathrm{N} ^ {\lambda + \mu} (\mathfrak {h})
\]

for all \(\lambda, \mu \in \mathrm{P}\) .

This follows from Props. 2 and 3.

PROPOSITION 9. Let \(\mathfrak{h}\) be a nilpotent Lie algebra and M a finite dimensional \(\mathfrak{h}\) -module. Denote by P the set of maps from \(\mathfrak{h}\) to \(k\) .

\begin{enumerate}
\def\labelenumi{(\roman{enumi})}
\item
  Each \(\mathrm{M}^{\lambda}(\mathfrak{h})\) is an \(\mathfrak{h}\) -submodule of \(\mathrm{M}\) . If \(x_{M}\) is triangularizable for all \(x\in \mathfrak{h}\) , then \(\mathrm{M} = \sum_{\lambda \in \mathbb{R}}\mathrm{M}^{\lambda}(\mathfrak{h})\) .
\item
  If \(k\) is infinite, there exists \(x \in \mathfrak{h}\) such that \(\mathrm{M}^0(x) = \mathrm{M}^0(\mathfrak{h})\) .
\item
  If \(k\) is of characteristic 0, and if \(\lambda \in \mathrm{P}\) is such that \(\mathrm{M}^{\lambda}(\mathfrak{h}) \neq 0\) , then \(\lambda\) is a linear form on \(\mathfrak{h}\) vanishing on \([\mathfrak{h},\mathfrak{h}]\) , and \(\mathrm{M}_{\lambda}(\mathfrak{h}) \neq 0\) .
\item
  If \(f: \mathrm{M} \to \mathrm{N}\) is a surjective homomorphism of finite dimensional \(\mathfrak{h}\) -modules, then \(f(\mathrm{M}^{\lambda}(\mathfrak{h})) = \mathrm{N}^{\lambda}(\mathfrak{h})\) for all \(\lambda \in \mathrm{P}\) .
\item
  If \(\mathrm{N}\) is a finite dimensional \(\mathfrak{h}\) -module, and \(\mathrm{B}\) a bilinear form on \(\mathrm{M} \times \mathrm{N}\) invariant under \(\mathfrak{h}\) , then \(\mathrm{M}^{\lambda}(\mathfrak{h})\) and \(\mathrm{N}^{\mu}(\mathfrak{h})\) are orthogonal relative to
\end{enumerate}

B if \(\lambda + \mu \neq 0\) . Moreover, if B is non-degenerate then so is its restriction to \(\mathrm{M}^{\lambda}(\mathfrak{h}) \times \mathrm{N}^{-\lambda}(\mathfrak{h})\) for all \(\lambda \in \mathrm{P}\) .

Part (i) follows from no. 1, Lemma 1 and Th. 1. Part (ii) follows from no. 2, Prop. 7. Part (iv) follows from no. 1, Cor. 3 of Th. 1. We prove (iii). We can assume that \(\mathrm{M} = \mathrm{M}^{\lambda}(\mathfrak{h})\) . Then, for all \(x \in \mathfrak{h}\) , \(\lambda(x) = (\dim \mathrm{M})^{-1}\mathrm{Tr}(x_{\mathrm{M}})\) ; this proves that \(\lambda\) is linear (which also follows from Prop. 5) and that \(\lambda\) vanishes on \([\mathfrak{h},\mathfrak{h}]\) . Consider the map \(\rho : \mathfrak{h} \to \mathrm{End}_k(\mathrm{M})\) defined by

\[
\rho (x) = x _ {\mathrm{M}} - \lambda (x) 1 _ {\mathrm{M}};
\]

from the above, this is a representation of \(\mathfrak{h}\) on M, and \(\rho(x)\) is nilpotent for all \(x \in \mathfrak{h}\) . By Engel's theorem (Chap. I, §4, no. 2, Th. 1), there exists \(m \neq 0\) in M such that \(\rho(x)m = 0\) for all \(x \in \mathfrak{h}\) , so \(m \in \mathrm{M}_{\lambda}(\mathfrak{h})\) .

The first assertion of (v) follows from no. 1, Prop. 2 (ii). To prove the second, we can assume that \(k\) is algebraically closed in view of Prop. 1 of no. 1; it then follows from the first and the fact that \(\mathrm{M} = \sum_{\lambda} \mathrm{M}^{\lambda}(\mathfrak{h})\) , \(\mathrm{N} = \sum_{\lambda} \mathrm{N}^{\lambda}(\mathfrak{h})\) , cf.~(i).

Remark. Assume that \(k\) is perfect and of characteristic 2. Let \(\mathfrak{h} = \mathfrak{sl}(2, k)\) , and let \(M\) be the \(\mathfrak{h}\) -module \(k^2\) (for the identity map of \(\mathfrak{h}\) ). If \(x = \begin{pmatrix} a & b \\ c & a \end{pmatrix}\) is an arbitrary element of \(\mathfrak{h}\) , denote by \(\lambda(x)\) the unique \(\lambda \in k\) such that \(\lambda^2 = a^2 + bc\) . A calculation shows immediately that \(M = M^\lambda(\mathfrak{h})\) ; on the other hand, \(M_\lambda(\mathfrak{h}) = 0\) and \(\lambda\) is neither linear nor zero on \([\mathfrak{h}, \mathfrak{h}]\) , even though \(\mathfrak{h}\) is nilpotent.

COROLLARY. Let \(\mathfrak{h}\) be a nilpotent Lie algebra, and M a finite dimensional \(\mathfrak{h}\) -module such that \(\mathrm{M}^0 (\mathfrak{h}) = 0\) . Let \(f:\mathfrak{h}\to \mathrm{M}\) be a linear map such that

\[
f ([ x, y ]) = x. f (y) - y. f (x) \quad \text { for } x, y \in \mathfrak {h}.
\]

There exists \(a \in \mathbb{M}\) such that \(f(x) = x.a\) for all \(x \in \mathfrak{h}\) .

Let \(N = M \times k\) . Make h operate on N by the formula

\[
x. (m, \lambda) = (x. m - \lambda f (x), 0).
\]

The identity satisfied by \(f\) implies that \(\mathbf{N}\) is an \(\mathfrak{h}\) -module (Chap. I, §1, no. 8, Example 2). The map \((m, \lambda) \mapsto \lambda\) from \(\mathbf{N}\) to \(k\) is a homomorphism from \(\mathbf{N}\) to the trivial \(\mathfrak{h}\) -module \(k\) . By Prop. 9 (iv), it follows that \(\mathrm{N}^0(\mathfrak{h})\) contains an element of the form \((a, 1)\) with \(a \in \mathrm{M}\) . In view of the hypothesis on \(\mathbf{M}\) ,

\[
(\mathrm{M} \times 0) \cap \mathrm{N} ^ {0} (\mathfrak {h}) = 0,
\]

so \(\mathrm{N}^0 (\mathfrak{h})\) is of dimension 1 and hence is annihilated by \(\mathfrak{h}\) . Thus, \(x.a - f(x) = 0\) for all \(x\in \mathfrak{h}\) , which proves the corollary.

PROPOSITION 10. Let \(\mathfrak{g}\) be a Lie algebra, \(\mathfrak{h}\) a nilpotent subalgebra of \(\mathfrak{g}\) . Denote by P the set of maps from \(\mathfrak{h}\) to \(k\) .

\begin{enumerate}
\def\labelenumi{(\roman{enumi})}
\item
  For \(\lambda, \mu \in \mathrm{P}\) , \([\mathfrak{g}^{\lambda}(\mathfrak{h}), \mathfrak{g}^{\mu}(\mathfrak{h})] \subset \mathfrak{g}^{\lambda + \mu}(\mathfrak{h})\) ; in particular, \(\mathfrak{g}^{0}(\mathfrak{h})\) is a Lie subalgebra of \(\mathfrak{g}\) containing \(\mathfrak{h}\) , and the \(\mathfrak{g}^{\lambda}(\mathfrak{h})\) are stable under \(\operatorname{ad} \mathfrak{g}^{0}(\mathfrak{h})\) . Moreover, \(\mathfrak{g}^{0}(\mathfrak{h})\) is its own normalizer in \(\mathfrak{g}\) .
\item
  If M is a g-module, \(\mathfrak{g}^{\lambda}(\mathfrak{h})\mathrm{M}^{\mu}(\mathfrak{h})\subset\mathrm{M}^{\lambda+\mu}(\mathfrak{h})\) for \(\lambda,\mu\in P\) ; in particular, each \(\mathrm{M}^{\lambda}(\mathfrak{h})\) is a \(\mathfrak{g}^{0}(\mathfrak{h})\) -module.
\item
  If B is a bilinear form on g invariant under h, \(\mathfrak{g}^{\lambda}(\mathfrak{h})\) and \(\mathfrak{g}^{\mu}(\mathfrak{h})\) are orthogonal relative to B for \(\lambda + \mu \neq 0\) . Assume that B is non-degenerate. Then, for all \(\lambda \in P\) , the restriction of B to \(\mathfrak{g}^{\lambda}(\mathfrak{h}) \times \mathfrak{g}^{-\lambda}(\mathfrak{h})\) is non-degenerate; in particular, the restriction of B to \(\mathfrak{g}^{0}(\mathfrak{h}) \times \mathfrak{g}^{0}(\mathfrak{h})\) is non-degenerate.
\item
  Assume that \(k\) is of characteristic 0. Then, if \(x \in \mathfrak{g}^{\lambda}(\mathfrak{h})\) with \(\lambda \neq 0\) , ad \(x\) is nilpotent.
\end{enumerate}

The map \((x,y)\mapsto [x,y]\) from \(\mathfrak{g}\times \mathfrak{g}\) to \(\mathfrak{g}\) is \(\mathfrak{g}\) -invariant by the Jacobi identity, hence \(\mathfrak{h}\) -invariant. The first part of (i) thus follows from Prop. 2 (ii). Part (ii) is proved similarly.

If \(x\) belongs to the normalizer of \(\mathfrak{g}^0(\mathfrak{h})\) in \(\mathfrak{g}\) , (ad \(y\) ). \(x = -[x, y] \in \mathfrak{g}^0(\mathfrak{h})\) for all \(y \in \mathfrak{h}\) , so \((\operatorname{ad} y)^n\) . \(x = 0\) for \(n\) sufficiently large. This proves that \(x \in \mathfrak{g}^0(\mathfrak{h})\) . Assertion (i) is now completely proved.

Assertion (iii) follows from Prop. 9 (v).

To prove (iv), we can assume that \(k\) is algebraically closed. Let \(x \in \mathfrak{g}^{\lambda}(\mathfrak{h})\) , with \(\lambda \neq 0\) . For all \(\mu \in \mathrm{P}\) and any integer \(n \geq 0\) , \((\operatorname{ad} x)^n \mathfrak{g}^\mu(\mathfrak{h}) \subset \mathfrak{g}^{\mu + n\lambda}(\mathfrak{h})\) ; let \(\mathrm{P}_1\) be the finite set of \(\mu \in \mathrm{P}\) such that \(\mathfrak{g}^\mu(\mathfrak{h}) \neq 0\) ; if \(k\) is of characteristic 0 and \(\lambda \neq 0\) , \((\mathrm{P}_1 + n\lambda) \cap \mathrm{P}_1 = \varnothing\) for \(n\) sufficiently large, so \((\operatorname{ad} x)^n = 0\) .

Lemma 2. Assume that k is of characteristic 0. Let g be a semi-simple Lie algebra over k, B the Killing form of g, m a subalgebra of g. Assume that the following conditions are satisfied:

\begin{enumerate}
\def\labelenumi{\arabic{enumi})}
\item
  the restriction of B to m is non-degenerate;
\item
  if \(x \in \mathfrak{m}\) , the semi-simple and nilpotent components\footnote{By Chap. I, §6, no. 3, Th. 3, every \(x \in \mathfrak{g}\) can be written uniquely as the sum of a semi-simple element \(s\) and a nilpotent element \(n\) that commute with each other; the element \(s\) (resp. \(n\)) is called the semi-simple (resp. nilpotent) component of \(x\).} of \(x\) in \(\mathfrak{g}\) belong to \(\mathfrak{m}\) .
\end{enumerate}

Then \(\mathfrak{m}\) is reductive in \(\mathfrak{g}\) (Chap. I, § 6, no. 6).

By Chap. I, §6, no. 4, Prop. 5 d), m is reductive. Let c be the centre of m. If \(x \in c\) is nilpotent, then x = 0; indeed, for all \(y \in m\) , ad x and ad y commute, their composition ad \(x \circ ad y\) is nilpotent, and \(\mathrm{B}(x, y) = 0\) , so x = 0. Now let x be an arbitrary element of c; let s and n be its semi-simple and nilpotent components. We have \(n \in m\) . Since ad n is of the form \(\mathrm{P}(\mathrm{ad}x)\) , where P is a polynomial with no constant term, (ad n).m = 0 and so \(n \in c\) , and then n = 0 by the above. Thus ad x is semi-simple. Consequently, the restriction to m of the adjoint representation of g is semi-simple (Chap. I, §6, no. 5, Th. 4 b)).

PROPOSITION 11. Assume that k is of characteristic 0. Let g be a semi-simple Lie algebra, h a nilpotent subalgebra of g. The algebra \(\mathfrak{g}^{0}(\mathfrak{h})\) satisfies conditions (1) and (2) of Lemma 2; it is reductive in g.

Let \(x, x' \in \mathfrak{g}\) , \(s\) and \(s'\) their semi-simple components, \(n\) and \(n'\) their nilpotent components. We have

\[
\begin{array}{r l} x ^ {\prime} \in \mathfrak {g} ^ {0} (x) & \Longleftrightarrow (\operatorname{ad} s) (x ^ {\prime}) = 0 \quad \text {(Prop. 4)} \\ & \Longleftrightarrow (\operatorname{ad} x ^ {\prime}) (s) = 0 \\ & \Longrightarrow (\operatorname{ad} s ^ {\prime}) (s) = 0 \\ & \Longleftrightarrow (\operatorname{ad} s) (s ^ {\prime}) = 0 \\ & \Longleftrightarrow s ^ {\prime} \in \mathfrak {g} ^ {0} (x) \quad \text {(Prop. 4)} \end{array}
\]

so \(x' \in \mathfrak{g}^{0}(x) \Rightarrow n' \in \mathfrak{g}^{0}(x)\) and (2) is proved. The Killing form of g is non-degenerate, so its restriction to \(\mathfrak{g}^{0}(\mathfrak{h})\) is non-degenerate (Prop. 10 (iii)). The fact that \(\mathfrak{g}^{0}(\mathfrak{h})\) is reductive in g thus follows from Lemma 2.

